\documentclass[11pt]{article}
\usepackage[T1]{fontenc}
\usepackage{lmodern,amsmath,amssymb,amsthm,mathtools,microtype}
\usepackage[margin=1in]{geometry}
\usepackage{xcolor}
\usepackage[colorlinks=true,allcolors=blue!45!black]{hyperref}
\hypersetup{pdftitle={Prime-factor shapes and Liouville signs on two rough neighbors}}
\newtheorem{theorem}{Theorem}[section]
\newtheorem{lemma}[theorem]{Lemma}
\newtheorem{proposition}[theorem]{Proposition}
\newtheorem{corollary}[theorem]{Corollary}
\theoremstyle{remark}\newtheorem{remark}[theorem]{Remark}
\newcommand{\one}{\mathbf 1}
\newcommand{\e}{\mathrm e}
\newcommand{\PP}{\mathbb P}
\newcommand{\E}{\mathbb E}
\newcommand{\calL}{\mathcal L}
\newcommand{\Dcal}{\mathcal D}
\newcommand{\cP}{\mathcal P}
\newcommand{\dist}{\mathbb D}
\allowdisplaybreaks[1]
\raggedbottom
\setlength\emergencystretch{1.8em}
\title{\bfseries Prime-factor shapes and Liouville signs\\on two rough neighbors}
\author{\small \href{mailto:admin@cosmosgeek.org}{admin@cosmosgeek.org}\qquad \href{mailto:b321014@iiit-bh.ac.in}{b321014@iiit-bh.ac.in}}
\date{\small Public research draft for specialist review}
\begin{document}
\maketitle
\begin{abstract}
We adapt the quantitative correlation argument of Tao and Ter\"av\"ainen
to a centered difference of two multiplicative functions that agree on
its carrier primes. The application conditions both integers $n,n+2$
to have no prime factor below $\exp((\log X)^\gamma)$, with a sufficiently
small fixed $\gamma>0$ and a cutoff frozen in each outer block.
Outside a small set of scales, finite prime-factor shape tests have their
ordinary independent laws, and the two actual Liouville signs are fair
and independent of these tests. The result includes shrinking windows
for a genuine near-full prime factor at each site. We prove the centered
progression estimate, spell out the modified graph and circle reduction,
and justify sampling at literal powers of two by endpoint stability.
The result concerns rare rough populations and nonunit residuals; it
does not give a twin-prime count or a statement at every scale.
\end{abstract}

\section{Frozen blocks, tests, and statements}
Fix a large real $X$, put $L=\log X$, and keep
\[
 z=\exp(L^\gamma),\qquad R(n)=\one_{P^-(n)>z}
\]
fixed throughout the outer block $\sqrt X\le N\le X$.
Here $R$ is completely multiplicative, with $R(1)=1$, and
$\lambda(n)=(-1)^{\Omega(n)}$ is the Liouville function, counting
prime factors with multiplicity. For large $X$, $R$ vanishes on even
integers. Define the actual reference population
\[
 S_z(N)=\sum_{N<n\le2N}R(n)R(n+2),\qquad
 V_2(z)=\frac12\prod_{2<p\le z}\left(1-\frac2p\right).
\]
The fundamental lemma gives, uniformly on this block and for every
fixed $A>0$,
\begin{equation}\label{eq:rough-main}
 S_z(N)=NV_2(z)+O_A(NL^{-A}),\qquad
 V_2(z)\sim\frac{2C_2e^{-2\gamma_E}}{(\log z)^2},
\end{equation}
provided $\gamma$ is a sufficiently small fixed positive number.
Here $C_2=\prod_{p>2}(1-(p-1)^{-2})$ and $\gamma_E$ is Euler's constant.
A proof of the uniform error used here is given below.

Fix $\alpha>0$ and an integer $B$. A test $u$ is a step function on
$[0,\infty)$, with values in $[0,1]$, which is one outside a union of
at most $B$ intervals in $[\alpha,2]$. Its endpoints and heights may
depend on $X$, but are prescribed independently of the integers being
counted. Put
\[
 U(p)=u(\log p/L)
\]
and extend $U$ completely multiplicatively. Define the finite factorial
integral
\begin{equation}\label{eq:D}
 D_u(t)=1+\sum_{j=1}^{\lfloor t/\alpha\rfloor}\frac1{j!}
 \int_{\substack{x_i\ge\alpha\ (i\le j)\\x_1+\cdots+x_j\le t}}
            \prod_{i=1}^j\frac{u(x_i)-1}{x_i}\,dx_i.
\end{equation}
Values on interval boundaries do not affect these integrals. We use
the same notation $v,V,D_v$ for a second test. The functions $D_u$ are
the ordinary one-integer prime-factor laws, not new distributions.

\begin{theorem}[Frozen tests]\label{th:frozen}
There are absolute numbers $r>0$ and $\gamma_0>0$ with the following
property. Let $0<\gamma<\gamma_0$ be fixed, and put $a=r/4$, $b=r/2$.
For each fixed finite family $\mathcal F$ of pairs of tests as above, there is a set
$\mathcal E_X\subset[\sqrt X,X]$ of normalized logarithmic measure
$O_{\alpha,B,|\mathcal F|}(L^{-b})$ such that, outside it, for all tests in the family
and $\varepsilon,\varepsilon'\in\{-1,1\}$,
\begin{align}
 &\sum_{N<n\le2N} R(n)R(n+2)U(n)V(n+2)
      \one_{\lambda(n)=\varepsilon}\one_{\lambda(n+2)=\varepsilon'}
       \nonumber\\
 &\hspace{2cm}=\frac14 D_u(\log N/L)D_v(\log N/L)S_z(N)
                   +O_{\alpha,B}(NL^{-a}).\label{eq:frozen}
\end{align}
The same assertion holds outside a set of $O_{\alpha,B,|\mathcal F|}(L^{1-b})$
literal powers $N=2^j$ in $[\sqrt X,X]$. The exceptional sets may
depend on the prescribed test family and on the outer block $X$.
All estimates are uniform in the endpoints and heights within that family.
\end{theorem}

Normalized logarithmic measure means $L^{-1}\int dN/N$. In particular
the discrete conclusion concerns a vanishing fraction of the
$\asymp L$ powers of two in the block. It is not inferred merely from
the continuous exceptional measure.

\begin{corollary}[Finite shape cells]\label{cor:shapes}
Condition $n\in(N,2N]$ on both $n,n+2$ being $z$-rough. Record at each
site the number of distinct prime factors in any prescribed fixed
finite list of intervals for $\log p/\log N$, bounded away from zero.
For all but a vanishing fraction of the powers of two in the outer
block, their joint law tends to the product of the two ordinary
one-integer laws described by \eqref{eq:D}. The two Liouville signs
are independent fair signs, independent of these finite cells.
The same holds after requiring both sites to be squarefree.
\end{corollary}

\begin{corollary}[Shrinking giant-prime windows]\label{cor:giants}
Reduce $r$ if necessary, and choose fixed
\[
       0<\gamma,\sigma<r/64,
\]
with $\gamma<\gamma_0$. Fix $0<A_i<B_i$ for $i=1,2$.
At site $n_i=n+2(i-1)$ require a genuine prime divisor $p_i$ with
\begin{equation}\label{eq:giant-window}
 \frac{\log(N/p_i)}{\log N}
       \in[A_i(\log N)^{-\sigma},B_i(\log N)^{-\sigma}].
\end{equation}
For all but a vanishing fraction of the powers of two in the block,
the proportion of the both-rough population satisfying both requirements is
\[
 (B_1-A_1)(B_2-A_2)(\log N)^{-2\sigma}(1+o(1)).
\]
Each of the four Liouville sign cells has one quarter of this mass,
with relative error $o(1)$. These prime divisors are unique for large
$X$, and their residual factors are
$\exp(\Theta((\log N)^{1-\sigma}))$, not units.
\end{corollary}

The cutoff $z$ remains frozen in all three statements. No law for an
ungridded moving cutoff $z(N)$, adaptive postselection of tests, or
simultaneous uncountable family is asserted.

\section{The centered progression input}
We first prove the quantitative fact needed in place of the
equidistribution premise of the printed correlation theorem.

\begin{lemma}\label{lem:one-site}
Uniformly in the test endpoints, $0\le D_u(t)\le1$ on compact
subintervals of $(0,2)$, and $D_u$ is Lipschitz there with constants
depending only on $\alpha,B$. The same Lipschitz control holds for
changes of the interval endpoints.
If $0<\gamma<1/25$, $X^{.35}\le M\le2X$, and $q\le L^{.01}$, then
for every residue $c\pmod q$,
\begin{equation}\label{eq:center-AP}
 \sum_{\substack{M<n\le2M\\n\equiv c\pmod q}}
       R(n)\bigl(U(n)-D_u(\log M/L)\bigr)
                    =O_{\alpha,B}(ML^{-.35}).
\end{equation}
\end{lemma}

\begin{proof}
Begin without the gate $R$. Except on integers divisible by a square
of a prime $\ge X^\alpha$, expand
\[
 U(n)=\prod_{p\mid n}\bigl(1+(U(p)-1)\bigr)
       =\sum_{\substack{d\mid n\\d\text{ squarefree}}}
                  \prod_{p\mid d}(U(p)-1).
\]
There are a bounded number of selected factors, depending on $\alpha$.
The exceptional square count on the full interval is
$O(MX^{-\alpha}+\sqrt M)$, with bounded multiplicity, and is negligible
in the estimates below. The same argument is valid down to $M\ge X^{.34}$.
All selected primes are coprime to $q$ for large $X$.

Restrict first to $d\le M\exp(-L^{.3})$. The complementary integer in
the progression has count $M/(qd)+O(1)$. Summing the floor errors costs
$O_{\alpha,B}(M\exp(-L^{.3}))$. In the complementary range write $n=dk$,
where $k\le2\exp(L^{.3})$. Isolate one selected prime $p$, and fix $k$
and the remaining selected product $d'$. On every nonempty interval,
$M/(kd')\gg X^\alpha$. The prime-counting upper bound is consequently
$O_\alpha(M/(kd'L))$. Summing $1/k$ gives $O(L^{.3})$, while the
reciprocal mass of the bounded number of other macro primes is
$O_\alpha(1)$. The entire tail therefore costs $O_{\alpha,B}(ML^{-.7})$,
even after discarding the progression restriction.

Prime harmonic partial summation on $[X^\alpha,4X]$ identifies the
remaining principal with \eqref{eq:D}. The omitted terminal simplex
strip has width $O(L^{-.7})$ and that same cost: one coordinate can
be integrated last, and all nonzero coordinates are at least $\alpha$.
The classical prime number theorem gives smaller errors uniformly at
the interval endpoints. Thus
\begin{equation}\label{eq:ordinary-AP}
 \sum_{\substack{M<n\le2M\\n\equiv c\pmod q}}U(n)
       =\frac MqD_u(\log M/L)+O_{\alpha,B}(ML^{-.7})
\end{equation}
when the progression is primitive to all selected macro primes; the
latter holds here for every $c$. Formula \eqref{eq:D}, or the ordinary
count with $q=1$, also shows $0\le D_u\le1$.
The finite integral has bounded boundary densities on each compact
$t$-range. Changing $t$ or any of the finitely many interval endpoints
changes it by at most a constant times the total change. This proves
the asserted Lipschitz bounds.

Now insert upper and lower fundamental-lemma weights for $R$, supported
on $e\le D=\exp(L^{.2})$, with $|\lambda_e^\pm|\le1$ and prime
factors at most $z$; see \cite[Theorem 19.1]{kouk}.
Since $U\ge0$, the sieve inequalities apply to $RU$ before centering.
Also $U(e)=1$, so $U(em)=U(m)$ exactly. In a primitive class modulo $q$,
only $(e,q)=1$ occurs. Apply \eqref{eq:ordinary-AP} at $M/e$; it remains
at least $X^{.34}$ eventually. The absolute harmonic sum of the sieve
weights is at most $1+\log D=O(L^{.2})$. Thus the summed count error
is $O(ML^{-.5})$. Replacing
$D_u(\log(M/e)/L)$ by $D_u(\log M/L)$ costs at most
$O(ML^{-.6})$ after that harmonic sum.

The two sieve principals are
\[
 \frac{M}{\varphi(q)}D_u(\log M/L)
       \prod_{p\le z}(1-1/p)
\]
up to an error smaller than every fixed inverse power of $L$ times $M$.
This uses the fundamental-lemma clock
$\log D/\log z=L^{.2-\gamma}\to\infty$ and the fact that every prime
dividing $q$ is below $z$. The identical calculation with $U=1$ gives
the same principal without $D_u$. Subtracting proves
\eqref{eq:center-AP}, with room in the displayed exponent.
In a nonprimitive class both rough terms are identically zero.
\end{proof}

The same fundamental lemma applied to the forbidden classes $0,-2$
for odd primes and the single class at two proves
\eqref{eq:rough-main}. For example take level $D=X^{.1}$.
The CRT remainder is $O(D\log D)$, uniformly for $N\ge\sqrt X$, and
the clock $.1L^{1-\gamma}$ tends to infinity. The error is smaller than
$NL^{-A}$ for any fixed $A$. The displayed asymptotic for $V_2(z)$ is
Mertens' formula and the absolutely convergent twin singular product.

\section{Short exponential sums for the centered difference}
Set $f_1=RU$, $f_0=R$, and, for a fixed scale $M$,
\[
 d_M=D_u(\log M/L),\qquad F_M=f_1-d_M f_0.
\]
Both components are real, completely multiplicative and $1$-bounded;
$F_M$ is $1$-bounded but need not be multiplicative.
We use the standard distance
\[
 \dist(f,g;Y)^2=\sum_{p\le Y}\frac{1-\Re(f(p)\overline{g(p)})}{p}
\]
for multiplicative functions bounded by one.

\begin{lemma}\label{lem:short}
Fix a sufficiently small absolute $d_0\in(0,10^{-3}]$ and let
$\calL=L^{d_0}$,
$H_* =\exp(L^{1/10})$,
$H_0=\exp(L^{(1-\epsilon_1)/10})$, with a sufficiently small fixed
$\epsilon_1\in(0,1/100)$.
There exists an absolute $\eta_0>0$ such that, uniformly for
$X^{.4}\le M\le2X$, $H_0\le h\le H_*$ and $\theta\in\mathbb R$,
\begin{equation}\label{eq:short}
 \frac1M\int_M^{2M}
       \left|\frac1h\sum_{x<n\le x+h}F_M(n)e(\theta n)\right|dx
                         \ll_{\alpha,B}L^{-\eta_0}.
\end{equation}
The estimate remains valid for the slightly shorter lengths
$h\calL^{-O(1)}$ used below. Here $e(t)=\exp(2\pi it)$.
\end{lemma}

\begin{proof}
We give the modification of the short-sum proof in
\cite[Sections 3.4--3.6]{tt}; its two components must not be treated
as equidistributed functions on their own. Put
$Q_0=\calL^{1/10}$ and $Q=h\calL^{-1/10}$.

On the minor arcs, the argument of \cite[Section 3.5]{tt} gives a
power of $\calL$ saving for the short exponential mean of any real
$1$-bounded multiplicative function. It uses no progression-mean
hypothesis. More explicitly, the Ramar\'e expansion uses
$P_1=\log^{100}h$, $P_2=h^{1/2}$; the omitted integers cost
$O(h\log P_1/\log P_2)$ and repeated primes cost $O(h/P_1)$.
On each prime interval $P<p\le P'\le2P$, H\"older reduces the remaining
mean to its fourth moment. The prime-tuple upper sieve and the
geometric-sum estimate on these minor arcs give a fixed power saving
in $\calL$, uniformly in the coefficients of modulus at most one.
This is exactly the bound before the progression hypothesis is first
used in \cite[Section 3.6]{tt}. Apply it separately to $f_1$ and $f_0$;
$|d_M|\le1$ preserves the estimate for their difference.

On a major arc write $\theta=a/q+\xi$ with
$q\le Q_0$, $(a,q)=1$, and $|\xi|\le1/(qQ)$.
Divide the short interval into pieces of length
$\widetilde h=h\calL^{-1/5}$. On each residue modulo $q$ the extra
phase varies by $O(\calL^{-1/10})$; the resulting mean error has this
same bound after summing the pieces. It suffices to estimate the
centered short means in each residue class.

Nonprimitive classes are zero because $R$ removes every prime factor
of $q$. In a primitive class use Dirichlet characters. For each
component $f_i\chi$, subtract its own long mean on $(M,2M]$.
If $\chi$ is real, the uniform Matom\"aki--Radziwi\l\l{} theorem
\cite[Theorem 1]{mr} gives an $L^1$ short-minus-long bound
$O(L^{-1/400})$. To check the numerical room, take its exceptional-set
parameter $\delta=L^{-1/400}$ and use
$\log\widetilde h\ge L^{.098}$: the exceptional fraction is
$O(L^{-.015})$ plus an exponentially small term, and its remaining
error $\log\log\widetilde h/\log\widetilde h$ is smaller.

If $\chi$ is nonreal, the distance argument of
\cite[Lemma C.1]{mrt}, with the small-modulus uniformity checked in
\cite[Section 3.6]{tt}, gives
\[
 \inf_{|t|\le M}\dist(f_i\chi,n^{it};M)^2
                         \ge\frac1{100}\log\log H_*.
\]
The corrected \cite[Theorem A.1]{mrt} therefore gives an $L^1$ short
mean $O(L^{-1/2500})$: its squared-mean bound has first term
$O(L^{-1/1000}\log L)$, with the other two terms smaller.
Averaging the short means gives the same bound for the long mean,
up to $O(\widetilde h/M)$, so their difference has that bound as well.
Character and residue summations cost at most $q^2\le L^{.0002}$;
a positive absolute saving remains.

Only after these component estimates do we combine the long means.
Their combination is the long character mean of $F_M$.
Equation \eqref{eq:center-AP}, summed over residues if necessary,
bounds it by $O_{\alpha,B}(qL^{-.35})$.
Thus it is negligible with the same room. Taking the minimum of the
minor-arc saving, $d_0/20$, and a smaller absolute major-arc saving
proves \eqref{eq:short}. None of this applies the printed
equidistribution premise to $f_1$ or $f_0$ separately.
\end{proof}

\section{The centered correlation adapter}
\begin{proposition}\label{prop:adapter}
For some absolute $r_0>0$, uniformly in each prescribed pair of tests,
\begin{equation}\label{eq:integrated}
 \frac1L\int_{\sqrt X}^X
 \left|\frac1N\sum_{N<n\le2N}
        F_N(n)R(n+2)V(n+2)\right|\frac{dN}{N}
                          \ll_{\alpha,B}L^{-r_0}.
\end{equation}
The reflected orientation has the same estimate.
\end{proposition}

\begin{proof}
Write $G=RV$. We reconstruct the marked graph reduction so that
the nonmultiplicative first input is explicit. Use the scales of
Lemma~\ref{lem:short}. Choose a small absolute $\kappa>0$ later and put
\[
 J=\lfloor\kappa\log\calL\rfloor,\qquad
 V_*=\frac{\epsilon_1\log\log H_*}{J},\qquad C=\exp(V_*/2).
\]
For large $X$, $V_*$ is bounded above and below by constants depending
on $\kappa,d_0,\epsilon_1$, and can be made as large as needed by
reducing $\kappa$. Let $\cP_i$ consist of the primes with
\[
 C^{2i-2}<\frac{\log p}{\log H_0}<C^{2i-1}\quad(1\le i\le J).
\]
These disjoint banks satisfy
$V_*/3\le v_i:=\sum_{p\in\cP_i}1/p\le V_*$ by Mertens, for large
$X$. Choose $C\ge10$. For every selection $p_i\in\cP_i$, its product
$P$ is at most $H_*$, since the geometric sum gives
\[
 \sum_i\log p_i\le
 \frac{\log H_*}{C-C^{-1}}\le\log H_*.
\]
Also $J\le\epsilon_1^2\log\log H_*$, as required by the decoupling
theorem, once $\kappa$ is sufficiently small.

Every carrier prime lies above $z$ and below $X^\alpha$.
Consequently $f_0(p)=f_1(p)=G(p)=1$, and complete multiplicativity gives
\begin{equation}\label{eq:carrier}
           F_N(Pn)=F_N(n),\qquad G(Pn)=G(n).
\end{equation}
This uses a unit dilation eigenvalue. It does not assert that
$F_N(p)=1$; in fact $F_N(p)=1-d_N$ may be zero. All banks remain
available, so no carrier-selection or small-mean fallback is used.
Furthermore, uniformly for $P\le H_*$,
\begin{equation}\label{eq:drift}
          \|F_N-F_{N/P}\|_\infty
                \ll_{\alpha,B}\frac{\log P}{L}\ll L^{-.9}.
\end{equation}

Consider the nonnegative marked quantity
\[
 S_{\emptyset}(N)=\sum_{p_i\in\cP_i\ (i\le J)}
 \left|\frac1N\sum_{N<n\le2N}
                   F_N(n)G(n+2P)\one_{P\mid n}\right|.
\]
Changing $n=Pm$ and using \eqref{eq:carrier} changes its inner sum
to $P^{-1}$ times the correlation at scale $N/P$, with the center
changed by \eqref{eq:drift}. Shifting the endpoints of the logarithmic
$N$ integral by $\log P$ costs $O(L^{-.9})$, since the normalized
correlation is bounded. Hence, if $\mathcal T$ denotes the left side
of \eqref{eq:integrated},
\begin{equation}\label{eq:marked-integral}
 \frac1L\int_{\sqrt X}^X S_{\emptyset}(N)\frac{dN}{N}
       =\left(\prod_{i=1}^Jv_i\right)
                       \bigl(\mathcal T+O_{\alpha,B}(L^{-.9})\bigr).
\end{equation}
All scales after dilation remain above $X^{.4}$ for large $X$.

Expand the divisibility product using
$\one_{p\mid n}=(\one_{p\mid n}-1/p)+1/p$.
Equivalently solve for the full divisibility term from the product of
all centered factors. The triangle inequality bounds $S_{\emptyset}$
by the all-centered term and, for each nonempty $I\subset[J]$, the term
with factor $1/P_I$ and retained divisibility $P_{I^c}\mid n$.
The all-centered term is $O(V_*^{.51J})$ by the bounded-input
decoupling inequality \cite[Theorem 3.3]{tt}. To use its shift $P$,
split $n$ modulo two; the resulting residues modulo each odd carrier
prime are fixed, and its two functions are bounded. The interval
length $\asymp N$ satisfies that theorem's size condition. Neither
function is required to be multiplicative in this application.

For a nonempty $I$, fix the complementary primes and dilate by
$P_{I^c}$. Equations \eqref{eq:carrier}--\eqref{eq:drift} leave,
apart from a harmless center error, a reciprocal-prime average
\begin{equation}\label{eq:circle-target}
 \sum_{d\in\Dcal_I}\frac1d
 \left|\frac1M\sum_{M<m\le2M}F_M(m)G(m+2d)\right|,
\end{equation}
where $M=N/P_{I^c}$ and $\Dcal_I=\{\prod_{i\in I}p_i:p_i\in\cP_i\}$.
We give the quantitative circle estimate for this expression.

For a dyadic part $T/2<d\le T$, attach signs or unit phases $c_d$
to realize the absolute values, and set
\[
 Q_T(\theta)=\sum_{\substack{d\in\Dcal_I\\T/2<d\le T}}
                     \frac{c_d}{d}e(2d\theta),\qquad
 v_T=\sum_{\substack{d\in\Dcal_I\\T/2<d\le T}}\frac1d.
\]
The carrier-product fourth moment \cite[Lemma C.2]{pilatte} gives
\begin{equation}\label{eq:Q4}
 \int_0^1|Q_T(\theta)|^4d\theta
                  \ll\frac{V_*^{4J}}{T\log^4T},\qquad
 \|Q_T\|_\infty\le v_T.
\end{equation}
The same bound holds with the $c_d$: expand the fourth moment and
bound each coefficient by one. Replacing $\theta$ by $2\theta$
preserves its integral. The prime banks above satisfy the separated
factor condition used in that lemma; earlier factors have total
logarithm at most a fixed small fraction of the last factor's logarithm.

Average a simultaneous translation of the two bounded functions over
$1\le t\le T$. Endpoint errors are $O(TV_*^J/M)$ after normalization.
Fourier inversion expresses the dyadic part as
$1/(MT)$ times an integral of $Q_T$ times the two short transforms.
Their $L^2$ norms are $O(\sqrt T)$, the second transform is bounded
by $O(T)$, and \eqref{eq:short} bounds the mean absolute value of
the first by $O(MT\delta)$, where $\delta=L^{-\eta_0}$.
Splitting at $|Q_T|=\epsilon$ and using \eqref{eq:Q4} yields the bound
\[
 O\left(\epsilon+
     \frac{\delta V_*^{4J}v_T}{\epsilon^4\log^4T}
     +\frac{TV_*^J}{M}\right).
\]
Choose $\epsilon=V_*^J\delta^{1/5}/\log T$.
After summing the dyads, use
$\sum_Tv_T\le V_*^{|I|}$ and
$\sum_T1/\log T=O(\log L)$ for $H_0\le T\le H_*$.
Thus \eqref{eq:circle-target} is at most
$O(V_*^J L^{-\eta_1})$ for some absolute $\eta_1>0$,
with exponentially small endpoint errors. This argument used only
boundedness, the short exponential estimate, and the stated
carrier-product fourth moment.

Summing the complementary reciprocal weights and all subsets $I$
therefore gives, for a fixed absolute $C_1$,
\[
 S_{\emptyset}(N)\ll V_*^{.51J}
          +(C_1V_*^2)^J L^{-\eta_1}
          +C_1^J V_*^J L^{-.9}.
\]
Divide \eqref{eq:marked-integral} by
$\prod v_i\ge(V_*/3)^J$.
Choose $\kappa$ so small that $3V_*^{-.49}<1$ uniformly for large
$X$, and that $\kappa d_0\log(3C_1V_*)<\eta_1/2$.
These choices are compatible because $V_*\asymp1/\kappa$ and
$\kappa\log(1/\kappa)\to0$. The resulting bound is a negative
fixed power of $L$. This proves \eqref{eq:integrated} for some
positive absolute $r_0$, without assigning it an unsupported numerical
value. The proof for shift $-2$ is identical, and changing a physical
endpoint by two is negligible.
\end{proof}

\section{Signed projections and the rough joint mean}
The actual signs require a second, directly applicable branch of the
published theorem. Write
\[
 M(f;Y,Q)=\inf_{\substack{q\le Q,\ \chi\ ({\rm mod}\ q)\\|t|\le Y}}
                   \sum_{p\le Y}\frac{1-\Re(f(p)\overline{\chi(p)}p^{-it})}{p}.
\]
The Liouville distance lower used in \cite[Remarks 3.2]{tt} gives
$M(\lambda;X^2,L^{1/125})\ge c_\lambda\log L-O(1)$ for an absolute
$c_\lambda>0$. Passing to $\lambda RU$ changes this squared distance
by at most
\[
 \sum_{p\le z}\frac1p+\sum_{X^\alpha\le p\le X^2}\frac1p
                    \le\gamma\log L+O_\alpha(1).
\]
Choose $\gamma<c_\lambda/4$ and choose the earlier $d_0$ also below
$c_\lambda/10$. Then \cite[Theorem 3.1(ii)]{tt} applies to a signed
axis $\lambda RU$ with the opposite input $RV$ or $\lambda RV$.
It gives an absolute power-log bound outside a power-log exceptional
set. Since the functions are bounded, it also gives an integrated
bound of the form \eqref{eq:integrated}. This use is a direct use of
the printed nonpretentious theorem, unlike the centered adapter.

For the unsigned mean, Proposition~\ref{prop:adapter} first replaces
$RU$ by $D_u(\log N/L)R$ against $RV$ at the other site.
Use its reflected version to replace the latter $RV$ by $D_vR$.
The two errors are integrated power-log errors because $0\le D_u\le1$.
We obtain the main $D_uD_v S_z(N)$.
Expand each sign indicator as $(1+\varepsilon\lambda)/2$; all three
signed projections have the just-proved integrated bound. Reduce
$r>0$ below the finitely many resulting exponents, and if necessary
below $1/100$. We have proved, for the discrepancy $E(N)$ in each
instance of \eqref{eq:frozen},
\begin{equation}\label{eq:final-integrated}
       \frac1L\int_{\sqrt X}^X |E(N)|\frac{dN}{N^2}\ll_{\alpha,B}L^{-r}.
\end{equation}
All estimates are uniform in fixed test complexity and endpoints.
At this point choose $\gamma_0$ smaller than
$1/25$, $c_\lambda/4$ and $r/64$. The preceding exponents are
available before this final choice: their unsigned estimates used only
$\gamma<1/25$, and the signed distance has a fixed margin.

\section{Continuous scales, literal powers of two, and endpoint grids}
Markov's inequality applied to \eqref{eq:final-integrated} with
threshold $L^{-a}$, $a=r/4$, gives continuous exceptional measure
$O(L^{-3r/4})$, which is stronger than the claimed $O(L^{-b})$,
$b=r/2$.

For the discrete assertion we use more than continuous measure.
For each frozen test, the normalized discrepancy $C(N)=E(N)/N$
obeys, when $N'$ is near $N$,
\begin{equation}\label{eq:endpoint-stability}
 |C(N')-C(N)|\ll_{\alpha,B}|N'/N-1|+1/N.
\end{equation}
The functions being summed are bounded and fixed in the block.
Changing the interval changes at most $O(|N'-N|+1)$ terms, and the
centers $D_u(\log N/L)$ and $D_v(\log N/L)$ are Lipschitz.
The same elementary bound holds for $S_z(N)/N$.

If $|C(2^j)|>L^{-a}$, equation \eqref{eq:endpoint-stability} gives a
one-sided or two-sided logarithmic neighborhood of length
$\gg_{\alpha,B}L^{-a}$ on which $|C(N)|\gg L^{-a}$.
For large $X$ these neighborhoods around different powers of two are
disjoint. Each contributes $\gg L^{-2a}$ to the unnormalized integral.
Equation \eqref{eq:final-integrated} therefore bounds their number by
$O(L^{1-r+2a})=O(L^{1-b})$. A finite union handles all tests and signs.
This proves Theorem~\ref{th:frozen}; a small continuous exceptional
set alone would not have justified this conclusion.

We now pass from $\log p/L$ to $\log p/\log N$.
Take a mesh $L^{-\eta}$ for $t=\log N/L\in[1/2,1]$.
At each mesh point freeze all the required test endpoints in the
outer logarithmic coordinate. There are $O(L^\eta)$ such tests,
not an independent grid for every endpoint. The source bounds above
are uniform in their endpoints. The exceptional fraction for their
union is consequently $O(L^{-b+\eta})$.

The discrepancy between a physical test and its frozen version is
supported where a macro prime lies in one of finitely many endpoint
strips of width $O(L^{-\eta})$. The ordinary strip mass is
$O_{\alpha,B}(L^{-\eta})$ by \eqref{eq:D}.
The unsigned two-site theorem, applied to $1$ minus the corresponding
forbidden-strip multiplicative test, bounds the actual conditional
strip mass by
\begin{equation}\label{eq:grid-error}
 O_{\alpha,B}(L^{-\eta}+L^{-a+2\gamma}).
\end{equation}
Every sign-cell indicator is nonnegative and at most one, so the same
positive strip bound controls all sign-cell approximation errors.

For finite factor-count cells first refine overlapping intervals into
disjoint intervals, then use generating functions with step
heights in $[0,1]$. Their degrees are bounded in terms of $\alpha$,
since each selected prime has size at least a fixed power of $N$.
Interpolation at a fixed finite set of distinct heights recovers all
prescribed count cells. Its coefficients are fixed, so the estimates
and their grid errors survive. Counts with and without multiplicity
differ only at macro prime squares, whose ambient count is power-small.
This proves Corollary~\ref{cor:shapes} when $\eta<b$ and $2\gamma<a$.
Squarefree removal under the rough gate costs
$O(N/z+\sqrt N)$, negligible even relative to the smaller populations
below.

For Corollary~\ref{cor:giants} use $\eta=r/16$ and
$0<\gamma,\sigma<r/64$. Then
\begin{equation}\label{eq:hierarchy}
       2(\gamma+\sigma)<a,\qquad 2\sigma<\eta<b.
\end{equation}
In outer coordinates the giant-prime interval endpoints are
$t-A_i t^{1-\sigma}L^{-\sigma}$ and
$t-B_i t^{1-\sigma}L^{-\sigma}$; these are uniformly Lipschitz
functions of $t\in[1/2,1]$. They remain bounded away from zero,
so the same one-dimensional mesh applies.

There is at most one prime in each giant window at a site. Use
one minus the no-prime-in-window test and inclusion--exclusion.
The ordinary one-site probability is exactly
\[
 m_i(N)=\int_{1-B_i(\log N)^{-\sigma}}^{1-A_i(\log N)^{-\sigma}}
                       \frac{dx}{x}
       \sim(B_i-A_i)(\log N)^{-\sigma};
\]
two selected primes cannot fit because both exponents exceed $1/2$.
The frozen theorem gives product $m_1m_2$, and a quarter of it in
each sign cell. The arithmetic and positive strip errors
\eqref{eq:grid-error} are $o((\log N)^{-2\sigma})$ by
\eqref{eq:hierarchy}, while the exceptional fraction still tends to
zero. Finally
$\log(n_i/p_i)=\log(N/p_i)+O(1)$ uniformly on the dyad, proving the
stated nonunit residual size. This completes the proof of both corollaries.

\section{Prior work and limitations}
Ter\"av\"ainen \cite[Theorem 1.11]{ter} proves logarithmic independence
of large-prime counts. Tao--Ter\"av\"ainen \cite[Theorem 1.8]{tt} gives
quantitative ordinary smooth-neighbor laws outside exceptional scales.
The ordinary factorial law, the decoupling inequality, the short-interval
theorems and the carrier-product fourth moment are all prior inputs.
The present contribution is the centered extension with simultaneous
rare rough conditioning, actual parity cells and the shrinking-window
consequence, including the discrete-scale sampling argument.

The cutoff is subpower and chosen after the positive error exponents.
The conclusions do not allow fixed polynomial roughness, every scale,
postselected tests, or unit residuals. In particular a near-full prime
factor times a growing residual is not a prime integer. No strict
twin-prime lower bound follows from this note.

A focused primary-source comparison did not locate this exact combination
of relative conditioning and quantitative scopes. This is not a global
priority claim. ChatGPT and AI review agents assisted with extraction, proof
reconstruction, literature comparison and writing. Such checks are
internal, not external peer review or formal verification.


\begin{thebibliography}{9}\small
\bibitem{tt} T. Tao and J. Ter\"av\"ainen,
\emph{Quantitative correlations and some problems on prime factors of consecutive integers},
arXiv:2512.01739v2 (2026).
\href{https://arxiv.org/html/2512.01739v2}{Primary text}.
\bibitem{ter} J. Ter\"av\"ainen,
\emph{On binary correlations of multiplicative functions},
Forum Math. Sigma \textbf{6} (2018), e10.
\href{https://arxiv.org/html/1710.01195v2}{Primary preprint}.
\bibitem{mr} K. Matom\"aki and M. Radziwi\l\l,
\emph{Multiplicative functions in short intervals},
Ann. of Math. \textbf{183} (2016), 1015--1056.
\href{https://arxiv.org/html/1501.04585v3}{Primary preprint}.
\bibitem{mrt} K. Matom\"aki, M. Radziwi\l\l{} and T. Tao,
\emph{An averaged form of Chowla's conjecture},
Algebra Number Theory \textbf{9} (2015), 2167--2196;
corrected arXiv version 1503.05121v3 (2022).
\href{https://arxiv.org/html/1503.05121v3}{Corrected primary text}.
\bibitem{pilatte} C. Pilatte,
\emph{Improved bounds for the two-point logarithmic Chowla conjecture},
arXiv:2310.19357, Appendix C.
\href{https://arxiv.org/html/2310.19357}{Primary text}.
\bibitem{kouk} D. Koukoulopoulos,
\emph{The Distribution of Prime Numbers}, American Mathematical Society, 2019.
\href{https://dms.umontreal.ca/~koukoulo/documents/publications/primes.pdf}{Author-hosted text}.
\end{thebibliography}
\end{document}
